
\documentclass[11pt,a4paper]{ctexart}

\usepackage[a4paper,margin=2.2cm]{geometry}
\usepackage{longtable}
\usepackage{booktabs}
\usepackage{array}
\usepackage{enumitem}
\usepackage{xcolor}
\usepackage{hyperref}
\usepackage{listings}
\usepackage{fancyhdr}
\usepackage{lastpage}
\usepackage{titlesec}
\usepackage{parskip}

\hypersetup{
  colorlinks=true,
  linkcolor=blue!55!black,
  urlcolor=blue!65!black,
  pdftitle={AgentBenchOps 项目交接与恢复手册},
  pdfauthor={AgentBenchOps}
}

\definecolor{codebg}{RGB}{247,248,250}
\definecolor{codeframe}{RGB}{220,223,228}
\definecolor{keywordcolor}{RGB}{48,86,145}
\definecolor{commentcolor}{RGB}{90,120,90}
\definecolor{stringcolor}{RGB}{150,70,70}

% listings does not include a built-in PowerShell or JSON lexer in all
% TeX Live distributions, so define them locally for portable compilation.
\lstdefinelanguage{PowerShell}{
  sensitive=false,
  morekeywords={
    begin,break,catch,class,continue,data,define,do,dynamicparam,else,
    elseif,end,exit,filter,finally,for,foreach,from,function,if,in,param,
    process,return,switch,throw,trap,try,until,using,var,while,workflow,
    Write-Host,Write-Output,Where-Object,ForEach-Object,Select-Object,
    Get-ChildItem,Get-Content,Set-Content,Copy-Item,New-Item,Remove-Item
  },
  morecomment=[l]{\#},
  morestring=[b]",
  morestring=[b]'
}

\lstdefinelanguage{JSON}{
  sensitive=true,
  morestring=[b]",
  morecomment=[l]{//},
  morecomment=[s]{/*}{*/},
  literate=
   *{0}{{{\color{keywordcolor}0}}}{1}
    {1}{{{\color{keywordcolor}1}}}{1}
    {2}{{{\color{keywordcolor}2}}}{1}
    {3}{{{\color{keywordcolor}3}}}{1}
    {4}{{{\color{keywordcolor}4}}}{1}
    {5}{{{\color{keywordcolor}5}}}{1}
    {6}{{{\color{keywordcolor}6}}}{1}
    {7}{{{\color{keywordcolor}7}}}{1}
    {8}{{{\color{keywordcolor}8}}}{1}
    {9}{{{\color{keywordcolor}9}}}{1}
    {:}{{{\color{black}:}}}{1}
    {,}{{{\color{black},}}}{1}
    {\{}{{{\color{black}\{}}}{1}
    {\}}{{{\color{black}\}}}}{1}
    {[}{{{\color{black}[}}}{1}
    {]}{{{\color{black}]}}}{1}
}

\lstdefinestyle{agentbench}{
  backgroundcolor=\color{codebg},
  frame=single,
  rulecolor=\color{codeframe},
  basicstyle=\ttfamily\small,
  breaklines=true,
  breakatwhitespace=false,
  columns=fullflexible,
  keepspaces=true,
  showstringspaces=false,
  tabsize=4,
  keywordstyle=\color{keywordcolor}\bfseries,
  commentstyle=\color{commentcolor},
  stringstyle=\color{stringcolor},
  numbers=left,
  numberstyle=\tiny\color{gray},
  xleftmargin=1.5em,
  framexleftmargin=1em
}
\lstset{style=agentbench}

\setlength{\headheight}{14pt}
\pagestyle{fancy}
\fancyhf{}
\lhead{AgentBenchOps}
\rhead{项目交接与恢复手册}
\cfoot{第 \thepage\ 页 / 共 \pageref{LastPage} 页}

\titleformat{\section}{\Large\bfseries}{\thesection}{0.8em}{}
\titleformat{\subsection}{\large\bfseries}{\thesubsection}{0.8em}{}

\title{\textbf{AgentBenchOps 项目交接与恢复手册}\\
\large Tool-Calling Agent 评测与可靠性平台}
\author{项目状态快照}
\date{2026-07-20}

\begin{document}

\maketitle
\tableofcontents
\newpage

\section{本文档用途}

本文档用于在聊天记录丢失、切换会话或更换协作者后，快速恢复
\texttt{AgentBenchOps} 项目的完整上下文。它记录：

\begin{itemize}[leftmargin=2em]
  \item 项目目标、技术栈和关键设计决策；
  \item 已完成阶段、涉及文件、执行命令和验收结果；
  \item 数据库表、Benchmark 任务、工具协议、幂等与恢复机制；
  \item 已出现并修复的问题；
  \item 当前项目进度、下一阶段以及后续路线图。
\end{itemize}

\textbf{重要原则：}代码仓库中的源代码和 Alembic 迁移是最终事实来源；
本文档是恢复上下文、解释设计和指导后续工作的交接材料。

\section{当前状态摘要}

\subsection{当前阶段}

截至 2026-07-20，项目已完成：

\begin{center}
\textbf{阶段 4E-1：扩展只读业务工具}
\end{center}

当前默认工具注册表包含：

\begin{lstlisting}
create_ticket
get_account
get_employee
get_ticket
list_employee_permissions
\end{lstlisting}

下一阶段是：

\begin{center}
\textbf{阶段 4E-2：update\_ticket（乐观锁、幂等更新和未知状态恢复）}
\end{center}

MVP 完成度约为 \textbf{55\%}。可靠性基础层已经建立，尚未正式接入 LangGraph Agent。

\subsection{最新验收结果}

最新执行命令：

\begin{lstlisting}[language=bash]
python scripts/demo_read_business_tools.py
\end{lstlisting}

验收结果：

\begin{itemize}[leftmargin=2em]
  \item \texttt{get\_account} 成功，返回 \texttt{account\_001}，用户名
        \texttt{zhangsan}，状态 \texttt{active}；
  \item \texttt{list\_employee\_permissions} 成功，员工存在且当前权限数量为 0；
  \item \texttt{create\_ticket} 成功，创建
        \texttt{ticket\_3b65e31088164810b8a102e0}；
  \item \texttt{get\_ticket} 成功，读取到同一个 ticket ID；
  \item 四次调用均生成独立的 \texttt{operation\_id}；
  \item 返回结果、权限检查、工具范围检查和数据库访问均正常。
\end{itemize}

因此，\textbf{阶段 4E-1 已通过}。

\section{项目目标与范围}

\subsection{总体目标}

构建一个面向 Tool-Calling Agent 的评测与运行平台，核心能力包括：

\begin{itemize}[leftmargin=2em]
  \item FastAPI 运行接口；
  \item LangGraph Agent 执行图；
  \item PostgreSQL 业务状态、运行记录、步骤、工具账本和评测结果；
  \item Redis 缓存、锁或后续恢复协调；
  \item MCP 兼容工具接入；
  \item OpenTelemetry / Langfuse 可观测性；
  \item Benchmark YAML 数据集和可重复状态重置；
  \item State Oracle 与 Trace Oracle；
  \item 权限、审批、故障注入、恢复、安全和统计实验；
  \item 最终至少 100 个 Benchmark 任务。
\end{itemize}

\subsection{MVP 业务域}

MVP 使用一个企业账号、权限和服务工单系统作为模拟业务域：

\begin{itemize}[leftmargin=2em]
  \item 员工 \texttt{employees}；
  \item 账号 \texttt{accounts}；
  \item 权限定义 \texttt{permissions}；
  \item 员工权限关系 \texttt{employee\_permissions}；
  \item 服务工单 \texttt{tickets}。
\end{itemize}

典型任务包括查询员工、查询账号、检查权限、创建工单、更新工单、
申请或撤销权限、账号恢复以及审批流程。

\section{开发环境}

\subsection{本地环境}

\begin{longtable}{>{\bfseries}p{4cm}p{9cm}}
\toprule
项目 & 当前配置 \\
\midrule
操作系统 & Windows，使用 VS Code；WSL2 Ubuntu 已安装 \\
Git & 2.51 \\
Python & 3.13.14 \\
虚拟环境 & 项目根目录下的 \texttt{.venv} \\
Docker & 29.6.1 \\
Docker Compose & v5.3 \\
PostgreSQL & \texttt{postgres:16-alpine} \\
Redis & \texttt{redis:7.4-alpine} \\
网络说明 & Docker Hub 拉取曾因直连 IPv6 失败，开启 Clash Verge 代理后恢复 \\
\bottomrule
\end{longtable}

\subsection{Python 依赖}

主要依赖：

\begin{lstlisting}
fastapi 0.139.2
sqlalchemy 2.0.51
asyncpg 0.31.0
redis 7.4.1
pydantic-settings
alembic
uvicorn
PyYAML
pytest 9.1.1
pytest-asyncio 1.4.0
ruff 0.15.22
httpx
\end{lstlisting}

项目约束：

\begin{lstlisting}
requires-python = ">=3.13,<3.14"
\end{lstlisting}

激活环境并安装：

\begin{lstlisting}[language=PowerShell]
.\.venv\Scripts\Activate.ps1
pip install -e ".[dev]"
\end{lstlisting}

测试必须优先使用：

\begin{lstlisting}[language=PowerShell]
python -m pytest
\end{lstlisting}

曾出现 \texttt{ModuleNotFoundError: app}，通过激活 \texttt{.venv}、
执行 editable install 并使用 \texttt{python -m pytest} 解决。

\section{项目目录结构}

当前主要目录：

\begin{lstlisting}
app/
  __init__.py
  config.py
  main.py

  agent/
    __init__.py
    graph.py
    nodes.py
    prompts.py
    state.py

  api/
    __init__.py
    approvals.py
    health.py
    runs.py
    tasks.py

  benchmark/
    __init__.py
    loader.py
    reset.py
    runner.py
    schemas.py

  domain/
    __init__.py
    accounts.py
    employees.py
    permissions.py
    tickets.py

  evaluation/
    __init__.py
    evaluator.py
    rules.py
    state_oracle.py
    trace_oracle.py

  observability/
    __init__.py
    events.py
    metrics.py
    tracing.py

  persistence/
    __init__.py
    base.py
    cache.py
    database.py
    models.py
    platform_models.py
    repositories.py

  tools/
    __init__.py
    gateway.py
    ledger.py
    recovery.py
    registry.py
    schemas.py
    implementations/
      __init__.py
      accounts.py
      employees.py
      permissions.py
      tickets.py

benchmark_tasks/
  basic/
    employee_lookup_001.yaml
    create_ticket_001.yaml
  permissions/
  recovery/
  security/

migrations/
scripts/
tests/
  test_health.py
  unit/
  integration/
  benchmark/
\end{lstlisting}

\section{阶段 0：架构与范围设计}

初始设计重点不是只让 Agent ``能调用工具''，而是确保运行结果可以被
验证、重放、恢复和解释。关键设计决策：

\begin{enumerate}[leftmargin=2.2em]
  \item 同时维护 \textbf{State Oracle} 和 \textbf{Trace Oracle}；
  \item 外部副作用必须有操作账本和明确的未知状态；
  \item 工具范围授权与业务权限必须分离；
  \item 工具调用必须具备参数模式、结果模式、风险等级、超时和审批元数据；
  \item Benchmark 必须版本化，初始状态必须可重复恢复；
  \item 评测指标必须精确定义，而不是只看最终文本；
  \item MVP 先跑通一个稳定业务域，再扩展到 100 个任务。
\end{enumerate}

\section{阶段 1：基础环境和容器}

\subsection{Docker Compose}

当前 Compose 仅包含 PostgreSQL 和 Redis：

\begin{itemize}[leftmargin=2em]
  \item PostgreSQL 用户、密码和数据库均为 \texttt{agentbench}；
  \item PostgreSQL 主机端口为 5432；
  \item Redis 主机端口为 6379；
  \item 两个服务都有 healthcheck；
  \item PostgreSQL 和 Redis 使用命名卷持久化。
\end{itemize}

常用命令：

\begin{lstlisting}[language=PowerShell]
docker compose up -d
docker compose ps
docker exec agentbench-postgres psql -U agentbench -d agentbench -c "SELECT 1;"
docker exec agentbench-redis redis-cli PING
\end{lstlisting}

预期数据库返回一行 \texttt{1}，Redis 返回 \texttt{PONG}。

\section{阶段 2：FastAPI 基础和健康检查}

\subsection{配置}

\texttt{app/config.py} 使用 Pydantic Settings，主要配置：

\begin{lstlisting}
app_name
env
debug
api_prefix
database_url
redis_url
sql_echo
\end{lstlisting}

本地 \texttt{.env} 使用宿主机地址访问 PostgreSQL 和 Redis。

\subsection{数据库和 Redis}

\texttt{app/persistence/database.py}：

\begin{itemize}[leftmargin=2em]
  \item 异步 SQLAlchemy Engine；
  \item \texttt{AsyncSessionFactory}；
  \item FastAPI 依赖 \texttt{get\_db\_session}。
\end{itemize}

\texttt{app/persistence/cache.py} 使用 \texttt{redis.asyncio.Redis.from\_url}。

\subsection{应用生命周期与健康接口}

\texttt{app/main.py} 使用 FastAPI lifespan：

\begin{itemize}[leftmargin=2em]
  \item 启动时创建 Redis 客户端；
  \item 关闭时关闭 Redis 和 SQLAlchemy Engine。
\end{itemize}

健康接口：

\begin{lstlisting}
GET /health/live
GET /health/ready
\end{lstlisting}

\texttt{/health/ready} 并发执行 PostgreSQL \texttt{SELECT 1} 和 Redis
\texttt{PING}，每项超时 3 秒：

\begin{itemize}[leftmargin=2em]
  \item 两者正常：HTTP 200，状态 \texttt{ready}；
  \item 任一失败：HTTP 503，状态 \texttt{not\_ready}；
  \item 同时返回各依赖的 \texttt{latency\_ms}。
\end{itemize}

已验证停止 Redis 后：

\begin{itemize}[leftmargin=2em]
  \item \texttt{/health/ready} 返回 503；
  \item \texttt{/health/live} 仍返回 200；
  \item 恢复 Redis 后 readiness 自动恢复。
\end{itemize}

\section{阶段 3A：业务数据模型}

\subsection{SQLAlchemy Base}

\texttt{app/persistence/base.py} 包含：

\begin{itemize}[leftmargin=2em]
  \item 统一约束命名 convention；
  \item \texttt{Base(DeclarativeBase)}；
  \item \texttt{TimestampMixin}，包含
        \texttt{created\_at} 和 \texttt{updated\_at}。
\end{itemize}

\subsection{业务表}

\begin{longtable}{p{3.3cm}p{10cm}}
\toprule
表 & 关键字段与约束 \\
\midrule
\texttt{employees} &
\texttt{id} 主键；\texttt{employee\_no} 唯一；
\texttt{name}；\texttt{department}；
状态为 \texttt{active/inactive/terminated}。 \\

\texttt{accounts} &
\texttt{id} 主键；\texttt{employee\_id} 唯一外键；
\texttt{username} 唯一；
状态为 \texttt{active/disabled/locked}；
\texttt{version > 0}。 \\

\texttt{permissions} &
\texttt{id}；\texttt{code} 唯一；
\texttt{name}；\texttt{description}；
风险等级 \texttt{low/medium/high/critical}；
\texttt{requires\_approval}。 \\

\texttt{employee\_permissions} &
员工与权限复合主键；
状态 \texttt{active/revoked}；
\texttt{granted\_by/granted\_at/revoked\_at}。 \\

\texttt{tickets} &
\texttt{id}；请求人与目标员工；
类型为 \texttt{permission\_grant/permission\_revoke/account\_recovery/general}；
状态、风险、标题、描述、解决结果和 \texttt{version}；
后续增加唯一 \texttt{source\_operation\_id}。 \\
\bottomrule
\end{longtable}

第一份迁移：

\begin{lstlisting}
d289cb2313aa create core business tables
\end{lstlisting}

已验证升级、降级到首迁移以及重新升级。

\section{阶段 3B：平台运行与评测数据模型}

平台表位于 \texttt{app/persistence/platform\_models.py}。

\subsection{benchmark\_tasks}

保存任务身份、数据集版本、用户请求、初始状态、可用工具、期望状态、
要求事件、禁止事件、时序规则、预算、元数据和激活状态。

唯一约束：

\begin{lstlisting}
(task_key, version)
\end{lstlisting}

\subsection{agent\_runs}

保存任务、实验、运行状态、模型、提供商、Prompt 版本、Agent 策略、
Memory 策略、输入、配置、随机种子、Checkpoint、Trace、步骤数、
工具调用数、Token、成本、延迟、最终回复、错误和时间字段。

\subsection{run\_steps}

保存运行步骤、父步骤、步骤序号、类型、状态、模型、工具、输入输出、
Token、延迟、错误和时间字段。运行内步骤序号唯一。

\subsection{tool\_operations}

工具操作账本，当前关键字段：

\begin{lstlisting}
id
operation_id
run_id
step_id
tool_name
arguments
arguments_hash
idempotency_key
risk_level
requires_approval
is_idempotent
status
retry_count
recovery_count
result
latency_ms
external_reference
error_type
error_message
error_details
recovered_at
recovery_details
started_at
finished_at
\end{lstlisting}

关键唯一约束：

\begin{lstlisting}
(run_id, idempotency_key)
operation_id
\end{lstlisting}

状态包含：

\begin{lstlisting}
prepared
running
succeeded
failed
unknown
rejected
cancelled
\end{lstlisting}

\subsection{approvals}

保存审批绑定、审批哈希、申请人与决策人、原因、过期与决策时间。

\subsection{evaluation\_results}

每个 run 一条评测结果，包含总通过状态、评测器版本、综合得分、
状态得分、Trace 得分、安全得分、违规项和最终状态。

第二份迁移：

\begin{lstlisting}
fff6ae231429
\end{lstlisting}

已执行升级、回滚到第一迁移、再次升级，并检查表结构。

后续迁移还增加：

\begin{itemize}[leftmargin=2em]
  \item \texttt{tool\_operations.latency\_ms}；
  \item \texttt{tool\_operations.error\_details}；
  \item \texttt{tickets.source\_operation\_id}；
  \item \texttt{tool\_operations.recovery\_count}；
  \item \texttt{tool\_operations.recovered\_at}；
  \item \texttt{tool\_operations.recovery\_details}。
\end{itemize}

\section{阶段 3C：Benchmark YAML 模式与导入}

\subsection{Pydantic 模式}

\texttt{app/benchmark/schemas.py} 目前包括：

\begin{lstlisting}
TaskBudget
EmployeeSeed
AccountSeed
PermissionSeed
EmployeePermissionSeed
TicketSeed
BusinessInitialState
BenchmarkTaskSpec
\end{lstlisting}

重要规则：

\begin{itemize}[leftmargin=2em]
  \item 所有模式使用 \texttt{extra="forbid"}；
  \item 状态字段使用 \texttt{Literal}；
  \item 检查重复 employee ID、员工号、username、permission code；
  \item 检查关系引用是否存在；
  \item \texttt{available\_tools} 会 trim，拒绝空名称与重复项；
  \item \texttt{initial\_state} 必须是
        \texttt{BusinessInitialState}。
\end{itemize}

\subsection{加载与导入}

\texttt{app/benchmark/loader.py}：

\begin{itemize}[leftmargin=2em]
  \item 使用 \texttt{yaml.safe\_load}；
  \item 递归加载 \texttt{.yaml/.yml}；
  \item 按路径排序；
  \item 拒绝重复的 \texttt{(task\_key, version)}。
\end{itemize}

Repository upsert 时必须序列化：

\begin{lstlisting}[language=Python]
"initial_state": spec.initial_state.model_dump(mode="json")
\end{lstlisting}

曾出现 Pydantic 对象不能写入 JSONB 的问题，已通过上述转换解决。

导入命令：

\begin{lstlisting}[language=PowerShell]
python scripts/import_benchmark_tasks.py
\end{lstlisting}

导入具有幂等性，多次执行不会重复创建任务。

\subsection{任务 employee\_lookup\_001}

文件：

\begin{lstlisting}
benchmark_tasks/basic/employee_lookup_001.yaml
\end{lstlisting}

任务内容：

\begin{itemize}[leftmargin=2em]
  \item 查询张三的员工号、部门和状态；
  \item 初始员工：
        \texttt{emp\_001 / E10001 / 张三 / 数据平台部 / active}；
  \item 初始账号：
        \texttt{account\_001 / zhangsan / active / version=1}；
  \item 可用工具：\texttt{get\_employee}；
  \item 禁止写工具；
  \item 预算：5 步、2 次工具、3000 Token、30 秒。
\end{itemize}

\subsection{任务 create\_ticket\_001}

文件：

\begin{lstlisting}
benchmark_tasks/basic/create_ticket_001.yaml
\end{lstlisting}

任务内容：

\begin{itemize}[leftmargin=2em]
  \item 张三为李四创建普通服务工单；
  \item 标题：数据平台访问问题；
  \item 描述：李四无法访问数据平台，请协助检查账号权限；
  \item 初始员工为 \texttt{emp\_001} 和 \texttt{emp\_002}；
  \item 初始工单为空；
  \item 可用工具：\texttt{get\_employee}、\texttt{create\_ticket}；
  \item 最终应存在状态为 \texttt{open} 的目标工单；
  \item 预算：8 步、4 次工具、5000 Token、45 秒。
\end{itemize}

\section{阶段 3D：业务状态重置}

\texttt{app/benchmark/reset.py} 提供：

\begin{lstlisting}
reset_business_state(session, state)
capture_business_state(session)
normalize_initial_state(state)
\end{lstlisting}

重置顺序：

\begin{lstlisting}
DELETE tickets
DELETE employee_permissions
DELETE accounts
DELETE permissions
DELETE employees

INSERT employees
INSERT permissions
FLUSH
INSERT accounts
INSERT employee_permissions
INSERT tickets
FLUSH
\end{lstlisting}

所有操作由调用方事务包裹，失败即回滚。

脚本：

\begin{lstlisting}[language=PowerShell]
python scripts/reset_business_state.py `
  --task-key employee_lookup_001 `
  --version 1
\end{lstlisting}

脚本会：

\begin{enumerate}[leftmargin=2.2em]
  \item 查找任务；
  \item 校验 \texttt{BusinessInitialState}；
  \item 重置业务表；
  \item 捕获最终状态；
  \item 与规范化初始状态比较；
  \item 打印实体数量。
\end{enumerate}

已验证重复重置和数据库结果一致。

\section{阶段 4A：工具协议、Registry 和 Gateway}

\subsection{统一工具协议}

\texttt{app/tools/schemas.py} 主要对象：

\begin{lstlisting}
ToolMetadata
ToolExecutionContext
ToolError
ToolExecutionResponse
ToolBusinessError
ToolDefinition
ToolHandler
\end{lstlisting}

当前 \texttt{ToolMetadata} 关键字段：

\begin{lstlisting}
name
description
risk_level
required_permissions
requires_approval
is_idempotent
read_only
timeout_seconds
external_reference_path
\end{lstlisting}

当前 \texttt{ToolExecutionContext} 关键字段：

\begin{lstlisting}
run_id
step_id
operation_id
actor_id
available_tools
permissions
fault_injection
\end{lstlisting}

授权被拆分为两层：

\begin{description}[leftmargin=3.7cm,style=nextline]
  \item[\texttt{available\_tools}]
        当前 Benchmark 任务允许调用的工具；
  \item[\texttt{permissions}]
        当前执行身份具备的业务权限。
\end{description}

两层都必须通过。

\subsection{第一个工具 get\_employee}

文件：

\begin{lstlisting}
app/tools/implementations/employees.py
\end{lstlisting}

支持且只能提供一个查询条件：

\begin{lstlisting}[language=JSON]
{"employee_id": "emp_001"}
{"employee_no": "E10001"}
{"name": "张三"}
\end{lstlisting}

错误码：

\begin{lstlisting}
employee_not_found
employee_ambiguous
invalid_arguments
\end{lstlisting}

元数据：

\begin{lstlisting}
required_permissions = {"employee.read"}
risk_level = low
read_only = true
is_idempotent = true
timeout_seconds = 3
\end{lstlisting}

\subsection{Registry}

\texttt{ToolRegistry}：

\begin{itemize}[leftmargin=2em]
  \item 注册工具定义；
  \item 禁止同名覆盖；
  \item 根据名称查询；
  \item 返回排序后的名称列表。
\end{itemize}

\subsection{Gateway 基础安全流程}

调用顺序：

\begin{lstlisting}
registered?
-> allowed by benchmark?
-> permissions satisfied?
-> arguments validate?
-> ledger claim
-> mark running
-> execute with timeout and transaction
-> result validate
-> ledger finalize
-> return structured response
\end{lstlisting}

未知工具、越权工具、缺少权限和错误参数都会被结构化拒绝。

曾修复的问题：

\begin{itemize}[leftmargin=2em]
  \item \texttt{employees.py} 被重复粘贴导致 E402；
  \item 删除重复内容时误删 \texttt{select}、\texttt{Employee} 和
        \texttt{ToolBusinessError} 导入；
  \item Python 3.13 + Ruff UP040 要求使用
        \texttt{type Alias = ...}，而不是
        \texttt{TypeAlias} 注解。
\end{itemize}

\section{阶段 4B：Tool Operation Ledger}

\subsection{目标}

每次带 \texttt{run\_id} 的工具调用都写入 \texttt{tool\_operations}，
记录参数、摘要、幂等键、状态、结果、错误与延迟。

\subsection{参数规范化和摘要}

\texttt{app/tools/ledger.py}：

\begin{lstlisting}[language=Python]
json.dumps(
    arguments,
    ensure_ascii=False,
    sort_keys=True,
    separators=(",", ":"),
)
\end{lstlisting}

之后计算 SHA-256，因此相同参数即使字典键顺序不同，也会得到相同摘要。

\subsection{幂等声明}

\texttt{claim()} 使用 PostgreSQL：

\begin{lstlisting}[language=SQL]
INSERT ... ON CONFLICT DO NOTHING
\end{lstlisting}

并依赖唯一约束：

\begin{lstlisting}
(run_id, idempotency_key)
\end{lstlisting}

行为：

\begin{itemize}[leftmargin=2em]
  \item 首次调用创建 \texttt{prepared} 记录；
  \item 执行前改为 \texttt{running}；
  \item 完成后改为最终状态；
  \item 重复请求增加 \texttt{retry\_count}；
  \item 同一幂等键和相同参数返回已有结果；
  \item 同一幂等键配不同工具或参数返回
        \texttt{idempotency\_key\_conflict}。
\end{itemize}

\subsection{已验证结果}

\begin{lstlisting}
调用两次
数据库只有一条 tool_operations
第二次 replayed = true
retry_count = 1
两个响应 operation_id 相同
\end{lstlisting}

\section{阶段 4C：第一个写工具 create\_ticket}

\subsection{事务边界}

Gateway 中写工具执行结构：

\begin{lstlisting}[language=Python]
async with asyncio.timeout(timeout_seconds):
    async with session.begin():
        raw_result = await handler(...)
        validated_result = result_model.model_validate(raw_result)
\end{lstlisting}

结果校验也在业务事务内部。返回结构错误时事务回滚。

\subsection{响应丢失故障注入}

故障模式：

\begin{lstlisting}
drop_response_after_commit
\end{lstlisting}

正确顺序：

\begin{lstlisting}
handler execution
-> result validation
-> transaction commit
-> injected TimeoutError
-> response status timed_out
-> ledger status unknown
\end{lstlisting}

最初错误实现把长时间 \texttt{sleep} 放在
\texttt{asyncio.timeout()} 外部，导致结果只是延迟约 5.1 秒后返回成功。
修复为事务提交后主动抛出 \texttt{TimeoutError}。

\subsection{unknown 判定}

写操作发生超时后：

\begin{lstlisting}[language=Python]
response.status == "timed_out"
and not definition.metadata.read_only
\end{lstlisting}

则账本必须写为 \texttt{unknown}，不能误写 \texttt{failed}。

原因：

\begin{itemize}[leftmargin=2em]
  \item \texttt{failed} 表示明确知道业务副作用没有成功；
  \item \texttt{unknown} 表示副作用可能已经提交，但调用方没收到结果；
  \item 是否幂等与结果是否未知是两个独立概念。
\end{itemize}

\section{阶段 4D：未知状态恢复}

\subsection{业务操作关联}

\texttt{tickets.source\_operation\_id} 唯一保存：

\begin{lstlisting}
tool_operations.operation_id
-> tickets.source_operation_id
\end{lstlisting}

Gateway 将账本生成的 \texttt{operation\_id} 注入
\texttt{ToolExecutionContext}，模型不能自行提供。

\texttt{create\_ticket} 创建工单时保存：

\begin{lstlisting}[language=Python]
source_operation_id=context.operation_id
\end{lstlisting}

如果同一 \texttt{operation\_id} 已经有工单，则直接返回已有工单，
避免重复副作用。

\subsection{恢复服务}

文件：

\begin{lstlisting}
app/tools/recovery.py
\end{lstlisting}

主要对象：

\begin{lstlisting}
RecoveryResolution
RecoveryResponse
RecoveryRegistry
UnknownOperationRecoveryService
recover_create_ticket
\end{lstlisting}

恢复流程：

\begin{enumerate}[leftmargin=2.2em]
  \item 根据 \texttt{operation\_id} 查询并锁定
        \texttt{ToolOperation}；
  \item 仅恢复状态为 \texttt{unknown} 的操作；
  \item 根据工具名查找恢复处理器；
  \item \texttt{recover\_create\_ticket} 使用
        \texttt{source\_operation\_id} 精确查找工单；
  \item 找到工单后把账本改为 \texttt{succeeded}；
  \item 保存结果、\texttt{external\_reference}、
        \texttt{recovery\_count}、\texttt{recovered\_at} 和
        \texttt{recovery\_details}；
  \item 相同幂等请求随后可以直接回放成功结果。
\end{enumerate}

\subsection{已验证恢复结果}

执行：

\begin{lstlisting}[language=PowerShell]
python scripts/demo_unknown_recovery.py
\end{lstlisting}

实际关键结果：

\begin{lstlisting}
first_response.status = timed_out
first_response.operation_id = op_339185010b0543d0850e35abc6e2ca46

recovery_response.previous_status = unknown
recovery_response.current_status = succeeded
recovery_response.recovered = true
recovery_response.recovery_count = 1

external_reference = ticket_6d7f9be621bf480abd7d2dfa
replay_response.status = succeeded
replay_response.replayed = true
ticket_count = 1
ledger.status = succeeded
ledger.retry_count = 1
ledger.recovery_count = 1
\end{lstlisting}

阶段 4D 完全通过。

\section{阶段 4E-1：扩展只读业务工具}

\subsection{get\_account}

文件：

\begin{lstlisting}
app/tools/implementations/accounts.py
\end{lstlisting}

查询条件二选一：

\begin{lstlisting}[language=JSON]
{"employee_id": "emp_001"}
{"username": "zhangsan"}
\end{lstlisting}

权限：

\begin{lstlisting}
account.read
\end{lstlisting}

错误码：

\begin{lstlisting}
account_not_found
account_ambiguous
invalid_arguments
\end{lstlisting}

\subsection{list\_employee\_permissions}

文件：

\begin{lstlisting}
app/tools/implementations/permissions.py
\end{lstlisting}

输入：

\begin{lstlisting}[language=JSON]
{
  "employee_id": "emp_001",
  "include_revoked": false
}
\end{lstlisting}

默认仅返回 active 权限。员工不存在返回
\texttt{employee\_not\_found}；员工存在但权限为空时成功返回
\texttt{count=0}。

权限：

\begin{lstlisting}
permission.read
\end{lstlisting}

\subsection{get\_ticket}

位于：

\begin{lstlisting}
app/tools/implementations/tickets.py
\end{lstlisting}

输入：

\begin{lstlisting}[language=JSON]
{"ticket_id": "ticket_..."}
\end{lstlisting}

权限：

\begin{lstlisting}
ticket.read
\end{lstlisting}

不存在时返回 \texttt{ticket\_not\_found}。

\subsection{当前默认工具注册表}

\begin{longtable}{p{4cm}p{3cm}p{2.3cm}p{2.5cm}}
\toprule
工具 & 所需权限 & 只读 & 风险 \\
\midrule
\texttt{get\_employee} & \texttt{employee.read} & 是 & low \\
\texttt{get\_account} & \texttt{account.read} & 是 & low \\
\texttt{list\_employee\_permissions} & \texttt{permission.read} & 是 & low \\
\texttt{get\_ticket} & \texttt{ticket.read} & 是 & low \\
\texttt{create\_ticket} & \texttt{ticket.write} & 否 & medium \\
\bottomrule
\end{longtable}

\subsection{阶段 4E-1 最新成功输出摘要}

\begin{lstlisting}
run_id = 7c60b971-dd5e-48cb-beb2-6ba8fd8effd7

get_account:
  status = succeeded
  account.id = account_001
  username = zhangsan
  account.status = active
  version = 1

list_employee_permissions:
  status = succeeded
  employee_id = emp_001
  count = 0

create_ticket:
  status = succeeded
  ticket.id = ticket_3b65e31088164810b8a102e0
  status = open
  risk_level = low

get_ticket:
  status = succeeded
  ticket.id = ticket_3b65e31088164810b8a102e0
\end{lstlisting}

\textbf{结论：}阶段 4E-1 通过。

\section{当前测试与质量检查}

常用完整检查：

\begin{lstlisting}[language=PowerShell]
ruff check app tests scripts --fix
ruff format app tests scripts
ruff check app tests scripts
python -m pytest
\end{lstlisting}

数据库迁移检查：

\begin{lstlisting}[language=PowerShell]
python -m alembic current
python -m alembic heads
python -m alembic upgrade head
\end{lstlisting}

重置和演示：

\begin{lstlisting}[language=PowerShell]
python scripts/import_benchmark_tasks.py

python scripts/reset_business_state.py `
  --task-key employee_lookup_001 `
  --version 1

python scripts/invoke_get_employee.py --name 张三
python scripts/demo_tool_ledger.py
python scripts/demo_create_ticket.py
python scripts/demo_unknown_recovery.py
python scripts/demo_read_business_tools.py
\end{lstlisting}

\section{关键设计不变量}

后续开发不得破坏以下约束：

\begin{enumerate}[leftmargin=2.2em]
  \item 工具只能通过 \texttt{ToolGateway} 调用；
  \item 工具定义必须有输入模式、输出模式和安全元数据；
  \item \texttt{available\_tools} 与 \texttt{permissions} 必须分别检查；
  \item 带 \texttt{run\_id} 的调用必须进入工具账本；
  \item 同一运行中的相同幂等键不得产生重复操作记录；
  \item 相同幂等键配不同参数必须拒绝；
  \item 写操作由 Gateway 管理事务；
  \item 结果模式校验失败时写事务必须回滚；
  \item 写操作提交后响应丢失必须记录为 \texttt{unknown}；
  \item unknown 恢复必须依据稳定业务标识，而不是模糊匹配标题或描述；
  \item \texttt{operation\_id} 由 Gateway 注入，不能信任模型输入；
  \item Benchmark 重置必须确定性、可重复并在事务中执行。
\end{enumerate}

\section{已知非阻塞优化项}

\subsection{回放延迟语义}

当前回放响应的 \texttt{latency\_ms} 使用原操作执行延迟。例如原调用为
27ms，幂等回放仍返回 27ms。后续可观测性阶段建议拆分：

\begin{lstlisting}
execution_latency_ms
replay_latency_ms
\end{lstlisting}

否则 P95/P99 可能把回放请求误计为原始执行耗时。

\subsection{运行统计}

部分演示脚本手动更新：

\begin{lstlisting}
AgentRun.total_tool_calls
AgentRun.status
AgentRun.finished_at
\end{lstlisting}

正式 Runner 阶段应由统一运行生命周期服务自动维护，避免脚本自行统计。

\section{下一阶段：4E-2 update\_ticket}

\subsection{目标}

实现第一个更新型写工具：

\begin{lstlisting}
update_ticket
\end{lstlisting}

必须覆盖：

\begin{enumerate}[leftmargin=2.2em]
  \item 根据 \texttt{ticket\_id} 查询工单；
  \item 要求调用方提供 \texttt{expected\_version}；
  \item 使用乐观锁防止并发覆盖；
  \item 至少支持更新状态、标题、描述或 resolution；
  \item 更新成功后 \texttt{version + 1}；
  \item 相同 operation ID 不重复执行；
  \item 响应丢失后状态为 \texttt{unknown}；
  \item 恢复器通过操作关联或更新审计判断是否已经成功；
  \item 恢复后可回放成功结果；
  \item 添加 Benchmark 任务和并发冲突测试。
\end{enumerate}

\subsection{建议输入模式}

\begin{lstlisting}[language=JSON]
{
  "ticket_id": "ticket_001",
  "expected_version": 1,
  "status": "in_progress",
  "resolution": null
}
\end{lstlisting}

必须至少提供一个可更新字段。

\subsection{建议错误码}

\begin{lstlisting}
ticket_not_found
ticket_version_conflict
no_ticket_changes
invalid_ticket_status_transition
\end{lstlisting}

\subsection{推荐实现策略}

使用原子条件更新：

\begin{lstlisting}[language=SQL]
UPDATE tickets
SET status = :status,
    version = version + 1,
    updated_at = now()
WHERE id = :ticket_id
  AND version = :expected_version
RETURNING ...
\end{lstlisting}

如果更新行数为 0，需要再次查询以区分：

\begin{itemize}[leftmargin=2em]
  \item 工单不存在；
  \item 工单存在但版本冲突。
\end{itemize}

为了支持 unknown 恢复，建议增加更新审计表，或为工单保存最近一次
\texttt{source\_operation\_id}。更稳健的方案是新建
\texttt{ticket\_mutations}：

\begin{lstlisting}
id
ticket_id
operation_id unique
previous_version
new_version
change_payload JSONB
created_at
\end{lstlisting}

恢复器可通过 \texttt{operation\_id} 精确确认更新是否已提交。

\section{后续路线图}

完成 4E-2 后，建议顺序：

\begin{enumerate}[leftmargin=2.2em]
  \item \textbf{LangGraph 最小执行闭环}\\
        Agent State、模型节点、工具节点、路由、终止条件；
  \item \textbf{Run/Step 生命周期}\\
        自动创建 AgentRun、RunStep、工具调用计数和最终状态；
  \item \textbf{Checkpoint 与恢复}\\
        Agent 中断后从最后稳定节点继续；
  \item \textbf{Evaluation Engine}\\
        State Oracle、Trace Oracle、预算和时序规则；
  \item \textbf{可观测性}\\
        统一事件模型、OpenTelemetry、Langfuse、指标；
  \item \textbf{审批与高风险权限工具}\\
        grant/revoke permission、审批绑定与防篡改；
  \item \textbf{故障和安全注入}\\
        超时、重复响应、工具错误、Prompt Injection、越权；
  \item \textbf{Benchmark 扩展}\\
        basic、permissions、recovery、security 四类任务扩展到 100 个；
  \item \textbf{统计实验}\\
        不同模型、Prompt、Agent 策略、Memory 策略和故障率对比。
\end{enumerate}

\section{新会话恢复操作清单}

当对话消失或重新开始时，将本文件提供给新的助手，并执行：

\begin{enumerate}[leftmargin=2.2em]
  \item 告知仓库路径：
        \texttt{E:\textbackslash Working\textbackslash agent-bench-ops}；
  \item 激活虚拟环境；
  \item 启动 Docker Compose；
  \item 查看 Git 状态和最近提交；
  \item 查看 Alembic 当前版本；
  \item 运行 Ruff 和 pytest；
  \item 导入 Benchmark；
  \item 运行最新演示脚本；
  \item 从 ``下一阶段：4E-2 update\_ticket'' 继续。
\end{enumerate}

命令：

\begin{lstlisting}[language=PowerShell]
cd E:\Working\agent-bench-ops
.\.venv\Scripts\Activate.ps1

docker compose up -d
docker compose ps

git status
git log --oneline -10

python -m alembic current
python -m alembic heads
python -m alembic upgrade head

ruff check app tests scripts
python -m pytest

python scripts/import_benchmark_tasks.py
python scripts/demo_read_business_tools.py
\end{lstlisting}

新助手需要重点确认：

\begin{lstlisting}
Current completed stage: 4E-1
Next stage: 4E-2 update_ticket
Default registered tools: 5
Unknown create_ticket recovery: passed
Latest read business tools demo: passed
\end{lstlisting}

\section{建议 Git 提交记录}

之前建议过的提交信息：

\begin{lstlisting}
feat: add core business schema
feat: add runtime and evaluation schema
feat: add benchmark state reset
feat: add tool operation ledger
feat: add transactional create ticket tool
feat: recover unknown tool operations
\end{lstlisting}

当前阶段可提交：

\begin{lstlisting}[language=PowerShell]
git add .
git commit -m "feat: add business read tools"
\end{lstlisting}

\section{LaTeX 编译说明}

本文件使用 \texttt{ctexart}，推荐使用 XeLaTeX：

\begin{lstlisting}[language=PowerShell]
xelatex agentbenchops_handoff.tex
xelatex agentbenchops_handoff.tex
\end{lstlisting}

运行两次是为了正确生成目录和页码引用。

\section{最终交接摘要}

\begin{center}
\fbox{
\begin{minipage}{0.9\textwidth}
\textbf{项目：}AgentBenchOps

\textbf{当前完成阶段：}4E-1 扩展只读业务工具

\textbf{最新验证：}get\_account、list\_employee\_permissions、
create\_ticket、get\_ticket 全部成功

\textbf{可靠性能力：}参数校验、工具范围、权限、事务、账本、参数哈希、
幂等回放、unknown 状态、业务副作用恢复

\textbf{下一阶段：}4E-2 update\_ticket，加入 expected\_version、
乐观锁、幂等更新和 unknown 恢复

\textbf{随后：}LangGraph 最小执行闭环
\end{minipage}
}
\end{center}

\end{document}
